\chapter{The Command System}
\label{ch:command-system}

\newthought{Everything in this part of the book eventually becomes a command.} A
calculator is a program you launch; a library dependency is a program you build; an
initialization step is a program you run once. Between the string you write in the card
and the process that actually starts, there is exactly one piece of machinery --- and
this chapter is about it, because every chapter that follows depends on knowing what it
does.

The problem it solves is stated quickly. A command string has to name things that live
in two different worlds:

\begin{description}[leftmargin=1.6em]
  \item[Things known when the card is read] where the project lives, where an external
        tool was installed, what the scan is called. These are the same for the whole
        run.
  \item[Things not known until a point exists] this Sample's identity, its working
        directory, which calculator slot the Worker happens to hold. These differ for
        every one of the million points you may be about to run.
\end{description}

Writing absolute paths would solve the first and ignore the second --- and would produce
a card nobody else can run. So \Jtwo\ resolves command strings in \textbf{two phases}, at
two different times, in two different processes. Section~\ref{sec:tokens} introduced that
idea; this chapter is the complete account.

\section{One parser, built once, shipped to the Workers}
\label{sec:parser-lifecycle}

There is a single object doing all of this --- a command parser, built once in the
control process when the card is loaded. Its construction is the moment every static
answer gets pinned down: the project root, the scan name, the path of every
\yamlkey{LibDeps} module, the absolute path of every registered executable.

Once built, the parser is \emph{serialized and handed to each Worker}. This matters more
than it looks. A Worker never re-reads your task card and never re-derives a path; it
receives a parser that already knows every static answer and only ever fills in the
per-Sample blanks. Two consequences you can rely on:

\begin{itemize}
  \item \textbf{A path resolves identically everywhere.} All Workers, on all machines,
        inherit one set of answers computed once. There is no way for two Workers to
        disagree about where a tool lives.
  \item \textbf{A broken path fails at load, not at minute forty.} Registering an
        executable checks that the source exists, so a typo in a tool path stops the run
        before a single Sample is submitted.
\end{itemize}

\section{Phase 1: resolved once, at load}
\label{sec:phase1}

Phase 1 runs in the control process, immediately after the card is validated. It walks
the configuration \emph{recursively} --- every string in every mapping and list, not just
the fields named \yamlkey{cmd} --- and substitutes every static token it finds. A path
buried in an I/O spec gets the same treatment as a command.
Table~\ref{tab:phase1-tokens} is the complete set.

\begin{table}[ht]
  \footnotesize
  \begingroup
  \setlength{\tabcolsep}{0pt}
  \begin{tabular}{@{}p{0.312\linewidth}p{0.688\linewidth}@{}}
    \toprule
    \textbf{Token} & \textbf{Resolves to} \\
    \midrule
    \marker{\&J/} (or bare \marker{\&J}) & the project root \\
    \code{\$\{Scan:name\}} & the scan name \\
    \code{\$\{Scan:root\}} & \file{<root>/\allowbreak{}outputs/\allowbreak{}<scan>};
      any other suffix (\code{output}, \code{outputs}, \ldots) resolves the same way \\
    \code{\$\{LibDeps:<name>\}} & the path of the tool registered under
      \yamlkey{LibDeps.Modules} (Chapter~\ref{ch:libdeps}). An unknown name is a hard
      error naming the reference \\
    \code{\$\{LibDeps:path\}} & the \yamlkey{LibDeps.path} base directory itself ---
      useful for reaching a sibling of your tools \\
    \emph{a registered executable name} & its absolute path
      (Section~\ref{sec:registered}) \\
    \code{@\{ROOT path\}} & the \textsc{cern} \textsc{root} installation prefix
      (Section~\ref{sec:root-path}) \\
    \jtt{\textasciitilde} & your home directory \\
    \code{\$\{source\}}, \code{\$\{path\}} & the module's \yamlkey{source} and runtime
      \yamlkey{path} --- \emph{only} inside \yamlkey{installation} and
      \yamlkey{initialization} commands (Chapter~\ref{ch:calculators}) \\
    \bottomrule
  \end{tabular}
  \endgroup
  \caption{Phase-1 tokens. All are resolved once, in the control process, before any
  Sample exists.}
  \label{tab:phase1-tokens}
\end{table}

\begin{jnote}
Phase 1 is also a checkpoint. If a static token survives its own phase --- a
\marker{\&J} that was never expanded, an unknown \code{\$\{LibDeps:\ldots\}} --- the load
stops there with the offending string quoted, rather than passing a half-resolved path to
a program that will fail more confusingly later.
\end{jnote}

\section{Phase 2: resolved per Sample, in the Worker}
\label{sec:phase2}

Phase 2 runs in the Worker, moments before the process starts, once per Sample. It knows
only three things, because those are the only three that vary per point ---
Table~\ref{tab:phase2-tokens} names them.

\begin{table}[ht]
  \footnotesize
  \begingroup
  \setlength{\tabcolsep}{0pt}
  \begin{tabular}{@{}p{0.229\linewidth}p{0.771\linewidth}@{}}
    \toprule
    \textbf{Token} & \textbf{Resolves to} \\
    \midrule
    \token{@SampleID} & the Sample's \textsc{uuid} \\
    \token{@Sdir} & the Sample's save directory under \file{SAMPLE/}.
      \textbf{Referring to it creates it} --- the directory is materialized on demand, so
      a Sample that never mentions \token{@Sdir} never costs an inode \\
    \token{@PackID} & the calculator slot the Worker currently holds --- a stable
      \code{001}\ldots\code{N}, and the isolation key for tools that cannot run twice in
      one directory (Chapter~\ref{ch:calculators}) \\
    \bottomrule
  \end{tabular}
  \endgroup
  \caption{Phase-2 tokens. Each is what keeps two Samples running the same program at the
  same instant from colliding.}
  \label{tab:phase2-tokens}
\end{table}

\section{The rule that connects them}
\label{sec:phase-rule}

One sentence governs the whole system:

\begin{quote}
Phase-2 tokens pass through Phase 1 untouched; a Phase-1 token reaching Phase 2 is an
error.
\end{quote}

The first half is what lets a single string carry both kinds, which is the normal case.
The second half is what stops a mistake from becoming a mystery. Phase 2 has no way to
resolve \code{\$\{LibDeps:SPheno\}} --- the Worker was never given the vocabulary --- so
rather than passing the literal text to a shell and letting the program fail with
something unhelpful, it refuses and says exactly which stage and field the bad string
came from:

\begin{terminal}
\begin{jout}
Phase 2 cannot resolve static tokens in stage 'execution'
field 'command': ${LibDeps:SPheno}/bin/SPheno @Sdir/in.slha
\end{jout}
\end{terminal}

In practice you meet this error in one situation: a string that was added to the card
somewhere Phase 1 does not walk, or that was assembled at run time out of pieces. The fix
is always to make sure the static half is written in the card, where Phase 1 can see it.

\subsection*{One string, both phases}

Because a string may mix them, Phase 1 splits it at the token boundaries, resolves the
static stretches, and leaves the per-Sample tokens standing. Written in the card:

\begin{terminal}[as written in the task card]
\begin{jcmd}
spheno ${LibDeps:SPheno}/templates/base.in @Sdir/@SampleID.out
\end{jcmd}
\end{terminal}

After Phase 1, in the control process --- note that \code{spheno} itself resolved too,
because it is a registered executable:

\begin{terminal}[after Phase 1, stored once for the whole run]
\begin{jcmd}
/scratch/tools/SPheno-4.0.5/bin/SPheno \
  /scratch/tools/SPheno-4.0.5/templates/base.in @Sdir/@SampleID.out
\end{jcmd}
\end{terminal}

And after Phase 2, in a Worker, for one particular point:

\begin{terminal}[what the Worker actually ran]
\begin{jcmd}
/scratch/tools/SPheno-4.0.5/bin/SPheno \
  /scratch/tools/SPheno-4.0.5/templates/base.in \
  /work/proj/outputs/scan/SAMPLE/000003/8f2c1d.../8f2c1d....out
\end{jcmd}
\end{terminal}

\section{Registered executables}
\label{sec:registered}

The tokens above are explicit --- you can see them in the string. Registered executables
are the one substitution that is \emph{implicit}: you declare a name once, and from then
on that bare word inside any command becomes an absolute path.

\begin{yamlwin}
LibDeps:
  path: "&J/External"
  Modules: [{name: SPheno}]
  registered_executables:
    - name: spheno
      source: "\${LibDeps:SPheno}/bin/SPheno"
      resolution: direct_path      # direct_path (default) | symlink
\end{yamlwin}

\begin{keylist}
  \item[name] \req. The bare word to substitute. It is replaced on \emph{word
        boundaries}, so a name never eats part of a longer word or a path segment ---
        registering \code{root} does not corrupt \file{/usr/local/rootfs}.
  \item[source] \req. Where the program actually is. It is itself Phase-1 resolved, so it
        may use \marker{\&J/} or \code{\$\{LibDeps:\ldots\}} as above. \textbf{The file
        must exist} --- registration fails at load if it does not.
  \item[resolution] \opt, default \code{direct\_path}. With \code{direct\_path} the name
        becomes the tool's own absolute path. With \code{symlink}, \Jtwo\ instead creates
        a link under \file{deps/registered/} inside the project and substitutes
        \emph{that}. Use \code{symlink} for tools that misbehave when invoked through a
        long or shifting path, or when you want the project to record which binaries it
        used.
\end{keylist}

\begin{jnote}
When several registered names overlap, the longest is substituted first. So registering
both \code{spheno} and \code{spheno\_hp} does what you meant: \code{spheno\_hp} is
matched as itself rather than as \code{spheno} followed by a stray suffix.
\end{jnote}

\section{\code{@\{ROOT path\}} and \yamlkey{EnvReqs.CERN\_ROOT}}
\label{sec:root-path}

Tools that link against \textsc{cern} \textsc{root} usually need its installation prefix
at build time. Hard-coding that prefix is exactly what makes a project unbuildable on the
next machine, so it gets a token of its own --- \code{@\{ROOT path\}} --- fed by an
\yamlkey{EnvReqs} block:

\begin{yamlwin}
EnvReqs:
  CERN_ROOT:
    required: true             # check that ROOT is installed
    version: "6.28"            # minimum version, optional
    get_path_command: "root-config --prefix"
\end{yamlwin}

\begin{keylist}
  \item[required] \opt, default \code{false}. When true, \Jtwo\ runs
        \cli{root-config --version} during preflight and fails the run if \textsc{root}
        is missing or older than \yamlkey{version}.
  \item[get\_path\_command] \opt. A command whose standard output \emph{is} the prefix.
        This is the portable form --- the machine tells you where its own \textsc{root}
        is, rather than the card asserting it.
  \item[path] \opt. A literal prefix, for when no such command exists. Note that this one
        is read \emph{even when} \yamlkey{required} is false, so a card may supply the
        prefix for a build without demanding a version check.
\end{keylist}

Using \code{@\{ROOT path\}} without configuring either \yamlkey{get\_path\_command} or
\yamlkey{path} is an error that says so directly.

\begin{jwarn}
\yamlkey{CERN\_ROOT} sits in the part of \yamlkey{EnvReqs} that \Jtwo\ passes through
rather than checks key by key (Chapter~\ref{ch:envreqs}), so a misspelled key inside it
is not reported --- it simply leaves the prefix unset, and you meet the problem as the
\code{@\{ROOT path\}} error above rather than as a validation message.
\end{jwarn}

\section{Where each token is allowed}
\label{sec:token-scope}

\begin{keylist}
  \item[Phase-1 tokens] any path-like or command field: \yamlkey{path},
        \yamlkey{source}, \yamlkey{cmd}, \yamlkey{cwd}, \yamlkey{data},
        \yamlkey{Sampling.Bounds.path}, \yamlkey{default\_yaml\_path}, and so on. Phase 1
        walks the whole configuration, so the practical answer is \emph{anywhere a string
        appears}.
  \item[\token{@SampleID}, \token{@Sdir}] calculator \yamlkey{execution} commands and I/O
        spec paths --- that is, anywhere that runs or writes once per point.
  \item[\token{@PackID}] module \yamlkey{path} and the commands of
        \yamlkey{clone\_shadow} modules, plus anywhere you want the slot recorded for
        traceability.
  \item[\code{\$\{source\}}, \code{\$\{path\}}] \yamlkey{installation} and
        \yamlkey{initialization} command lists only. These stages run before any Sample
        exists, which is also why Phase-2 tokens are rejected there.
\end{keylist}

\begin{jtip}
If you are unsure whether a string resolved the way you expected, run
\cli{Jarvis check} on the card. The smoke path executes the real commands for a couple
of points (Chapter~\ref{ch:check-modules}) and mirrors them to your terminal, which is
the fastest way to see a fully resolved command line without committing to a scan.
\end{jtip}
